\documentclass[11pt]{article}
\usepackage[a4paper,margin=2.1cm,top=1.9cm,bottom=1.9cm]{geometry}
\usepackage{booktabs}
\usepackage{tabularx}
\usepackage{array}
\usepackage{xcolor}
\usepackage{titlesec}
\usepackage[hidelinks]{hyperref}

\pagestyle{empty}
\setlength{\parindent}{0pt}
\renewcommand{\arraystretch}{1.35}

\newcolumntype{L}{>{\raggedright\arraybackslash}p{3.6cm}}
\newcolumntype{R}{>{\raggedright\arraybackslash}X}

\begin{document}

\begin{center}
    {\LARGE\bfseries Group Member Roles and Responsibilities}\\[6pt]
    {\large Generative Artificial Intelligence --- Final Project}\\[3pt]
    {\normalsize Korea University}\\[8pt]
    {\itshape GLOW vs FLUX: Comparing Normalizing Flows and Flow Matching\\
    for Face Image Generation at $64\times64$ Resolution}
\end{center}

\vspace{4pt}
\hrule
\vspace{10pt}

\textbf{Team composition.}\quad
This final project was carried out as an \textbf{individual (one-person) project}.
All roles and responsibilities listed below were performed by the single member.

\vspace{8pt}

\begin{tabularx}{\textwidth}{@{}lR@{}}
    \toprule
    \textbf{Member} & \textbf{Affiliation / Contact} \\
    \midrule
    Younghyeok Kim & Korea University \quad\textbullet\quad \texttt{younghyeok25@gmail.com} \\
    \bottomrule
\end{tabularx}

\vspace{14pt}

{\large\bfseries Roles and Responsibilities}

\vspace{6pt}

\begin{tabularx}{\textwidth}{@{}LR@{}}
    \toprule
    \textbf{Role} & \textbf{Responsibilities and Deliverables} \\
    \midrule
    Research Design \& Project Lead &
    Defined the research question (normalizing flows vs.\ flow matching at low
    resolution), designed the experimental protocol across four evaluation axes
    (sample quality, exact likelihood, reconstruction fidelity, sampling
    efficiency), and managed the overall project schedule. \\
    \addlinespace[2pt]
    Model Implementation &
    Implemented GLOW from scratch in PyTorch ($\sim$76M parameters): ActNorm,
    invertible $1\times1$ convolution with LU decomposition, affine coupling,
    and the multi-scale block architecture, plus the FFHQ-64 data pipeline
    (\texttt{model.py}, \texttt{dataset.py}). \\
    \addlinespace[2pt]
    Training \& Infrastructure &
    Built the training loop (NLL objective, 1{,}000-step LR warmup, batch 128,
    30{,}000 iterations) and resolved divergence via two stability fixes:
    log-scale clamping to $[-3,3]$ and a NaN/Inf gradient guard. Set up
    checkpointing and a watchdog script for auto-restart on crash
    (\texttt{train.py}, \texttt{run\_train\_with\_watchdog.sh}). \\
    \addlinespace[2pt]
    Baseline Sampling \& Evaluation &
    Ran pretrained baselines at $64\times64$ --- FLUX.1-schnell/dev (4--28 NFE)
    and DDIM-10/20/50 / DDPM-100 --- and measured FID (clean-fid, 5k samples vs.\
    70k FFHQ references), exact NLL, reconstruction PSNR/SSIM, and wall-clock
    sampling cost (\texttt{sample\_flux.py}, \texttt{sample\_ddim.py},
    \texttt{evaluate.py}, \texttt{reconstruct.py}). \\
    \addlinespace[2pt]
    Analysis \& Visualization &
    Conducted latent-space interpolation, OOD detection via NLL, and 2D
    two-moons transport-path comparison of the three paradigms; produced all
    paper figures 1--9 (\texttt{interpolation.py}, \texttt{ood\_detection.py},
    \texttt{visualize\_2d.py}, \texttt{make\_figures.py}). \\
    \addlinespace[2pt]
    Writing \& Presentation &
    Wrote the final report (\texttt{report.pdf}, two-column paper format),
    prepared the presentation slides and script, and authored the code
    documentation and reproduction guide accompanying the submission. \\
    \bottomrule
\end{tabularx}

\vspace{14pt}
\hrule
\vspace{6pt}
{\small All implementation, experiments, analysis, and writing were performed
solely by Younghyeok Kim.}

\end{document}
